\documentclass[11pt,a4paper]{article}
\usepackage{amsmath,amssymb,amsthm}
\usepackage{mathtools}
\usepackage{bm}
\usepackage{booktabs}
\usepackage{array}
\usepackage{enumitem}
\usepackage{geometry}
\usepackage{xcolor}
\usepackage{tcolorbox}
\usepackage{algorithm}
\usepackage{algpseudocode}
\usepackage[pdfencoding=auto,psdextra]{hyperref}

\geometry{margin=2.5cm}

\newcommand{\E}{\mathbb{E}}
\newcommand{\R}{\mathbb{R}}
\newcommand{\N}{\mathcal{N}}
\newcommand{\KL}{\mathrm{KL}}
\newcommand{\norm}[1]{\left\lVert #1 \right\rVert}
\newcommand{\inner}[2]{\left\langle #1,\,#2 \right\rangle}
\newcommand{\sg}{\operatorname{sg}}
\newcommand{\dd}{\mathrm{d}}
\newcommand{\Sbase}{S_0}
\newcommand{\Srl}{S_R}
\newcommand{\Lbase}{L_0}
\newcommand{\Lactor}{L_\theta}
\newcommand{\Lbeh}{L_{\bar\theta}}
\newcommand{\dmuS}{\Delta\mu_S}
\newcommand{\dvS}{\Delta v_S}
\newcommand{\dbS}{\Delta b_S}
\DeclareMathOperator*{\argmaxop}{arg\,max}

\theoremstyle{plain}
\newtheorem{theorem}{Theorem}[section]
\newtheorem{proposition}[theorem]{Proposition}
\newtheorem{lemma}[theorem]{Lemma}
\newtheorem{corollary}[theorem]{Corollary}

\theoremstyle{definition}
\newtheorem{definition}[theorem]{Definition}
\newtheorem{remark}[theorem]{Remark}

\title{\textbf{Weak-to-Strong Flow Direct-OPD}\\[0.3em]
\large Transferring a small model's RL-induced policy shift to a larger flow model}
\author{Jayce-Ping}
\date{2026-07-30}

\begin{document}
\maketitle

\begin{abstract}
Suppose reinforcement learning is affordable for a small flow model but expensive for a much
larger model.  Let $\Sbase$ be the small base checkpoint, $\Srl$ the same model after RL, and
$\Lbase$ a larger base model that we wish to improve.  Directly distilling $\Srl$ into the large
model is undesirable: it transfers both the useful RL change and the small model's capacity
limitations.  Motivated by Direct On-Policy Distillation for language models, we instead transfer
the \emph{policy shift} from $\Sbase$ to $\Srl$ and apply it on states visited by the large
recipient.

The formulation does not require flow models themselves to be Gaussian.  At the most general
level, the donor signal is the Radon--Nikodym ratio
$\dd K_{\Srl}/\dd K_{\Sbase}$ between two conditional Markov kernels, and the exact
KL-regularized target is the geometric tilt
\[
K_{\Lbase}
\left(\frac{\dd K_{\Srl}}{\dd K_{\Sbase}}\right)^\lambda.
\]
Gaussianity enters only when a stochastic sampler is chosen to make this target tractable.  For
continuous SDEs with shared diffusion, Girsanov's theorem turns the local donor shift into a drift
correction $b_{\Lbase}+\lambda(b_{\Srl}-b_{\Sbase})$; the globally tilted path law generally
requires an additional Doob continuation-value correction.  For Euler SDE and SDE-DPM-Solver++
kernels, the conditional target is Gaussian and the mean shift is available in closed form.
For deterministic ODE sampling, transition likelihood ratios are singular, but the same
velocity-delta target survives as a rescaled small-noise control-energy limit.  We derive all
four levels, give a recipient-on-policy training algorithm, and state compatibility checks and
experiments for FLUX.2-klein-base-9B $\rightarrow$ FLUX.2-dev.
\end{abstract}

\tableofcontents
\newpage

%% ============================================================
\section{Motivation, roles, and literature}
\label{sec:motivation}
%% ============================================================

\subsection{The weak-to-strong setting}

We use role-based notation rather than naming models by size:
\begin{center}
\begin{tabular}{@{}cll@{}}
\toprule
symbol & role & concrete example\\
\midrule
$\Sbase$ & small donor reference before RL & FLUX.2-klein-base-9B\\
$\Srl$ & same donor after task RL & a GenEval2 Flow-DPPO 9B LoRA\\
$\Lbase$ & large recipient reference & FLUX.2-dev, approximately 32B\\
$\Lactor$ & trainable large recipient & initialized from $\Lbase$\\
$\Lbeh$ & frozen rollout behavior checkpoint & refreshed from $\Lactor$\\
\bottomrule
\end{tabular}
\end{center}

The small pair is a \emph{donor} and the large model is a \emph{recipient}.  Calling the 9B
checkpoint ``teacher'' and the 32B checkpoint ``student'' obscures the central point: the final
9B policy need not be better than the 32B model.  The transferable object is the change induced
by RL,
\begin{equation}
\dvS(t,x,c)
:=
v_{\Srl}(t,x,c)-v_{\Sbase}(t,x,c),
\label{eq:velocity-shift}
\end{equation}
not the absolute field $v_{\Srl}$.

Directly matching $v_{\Lactor}$ to $v_{\Srl}$ imports the donor's pretraining biases, capacity
ceiling, and any missing knowledge.  Base subtraction removes what the donor already preferred
before RL and isolates what RL changed.

\subsection{The matching language-model literature}

The same idea has appeared in increasingly direct forms:
\begin{itemize}[leftmargin=1.6em]
\item Proxy-Tuning~\cite{proxytuning} adds
      $\text{logits}_{\mathrm{small,tuned}}-\text{logits}_{\mathrm{small,base}}$
      to a larger base model at decoding time.
\item RAST~\cite{rast} specializes the donor change to an RL-induced reasoning shift and tests
      transfer across model scales.
\item THINKLOGIT~\cite{thinklogit} uses the same arithmetic for long-reasoning transfer and
      optionally trains the small guider against target-aware preferences.
\item Logit Grafting~\cite{logitgrafting} shows that additive post-training logit deltas are
      exponential tilts and the exact solution of a KL-regularized reward objective.
\item Direct-OPD~\cite{directopd}, posted to arXiv in July 2026, is the closest match to the present
      goal.  It treats
      $\log\pi_{\mathrm{small,RL}}-\log\pi_{\mathrm{small,ref}}$
      as a dense implicit reward and trains a stronger student on its own on-policy states.
\end{itemize}

Categorical logits are natural parameters of the next-token distribution, so adding a logit
delta gives an exact exponential tilt.  A flow velocity is not generally a log density or a
natural parameter.  The remainder of this note identifies the conditions under which velocity
arithmetic has the same exact interpretation and what remains when those conditions do not hold.

\begin{tcolorbox}[colback=gray!5,colframe=gray!55,title=Scope]
The proposed method transfers an observed RL-induced policy shift.  Calling that shift the
``original RL reward'' requires the additional approximation that $\Srl$ is near the optimum of a
KL-regularized RL problem anchored at $\Sbase$.  GRPO clipping, reward normalization, finite
optimization, or zero-KL RL weaken that literal reward interpretation, but do not invalidate the
checkpoint-pair policy shift itself.
\end{tcolorbox}

%% ============================================================
\section{General non-Gaussian kernel formulation}
\label{sec:kernel}
%% ============================================================

\subsection{Recipient-on-policy query states}

Let
\begin{equation}
z=(t,x,c)
\label{eq:query-state}
\end{equation}
denote a physical time, latent state, and raw conditioning input.  At outer iteration $k$, we
sample $z$ from a frozen recipient behavior process $\Lbeh$.  All donor and recipient models are
then queried on the same physical $(t,x,c)$, although each architecture may encode the raw
condition $c$ with its own text encoder.

For the moment, let
\begin{equation}
K_j(\dd y\mid z),
\qquad
j\in\{\Sbase,\Srl,\Lbase,\Lactor\},
\label{eq:general-kernels}
\end{equation}
be arbitrary conditional Markov kernels on the same next-state space.  They need not be Gaussian.
Assume
\begin{equation}
K_{\Srl}(\cdot\mid z)\ll K_{\Sbase}(\cdot\mid z),
\qquad
K_{\Lbase}(\cdot\mid z)\ll K_{\Sbase}(\cdot\mid z),
\label{eq:absolute-continuity}
\end{equation}
so a version of the donor ratio is defined
$K_{\Lbase}(\cdot\mid z)$-almost everywhere.

\begin{definition}[Donor policy shift]
\label{def:donor-ratio}
The small-model RL shift at query $z$ is the Radon--Nikodym derivative
\begin{equation}
R_S(y\mid z)
:=
\frac{
  \dd K_{\Srl}(\cdot\mid z)
}{
  \dd K_{\Sbase}(\cdot\mid z)
}(y),
\qquad
r_S(y\mid z):=\log R_S(y\mid z).
\label{eq:donor-ratio}
\end{equation}
\end{definition}

The ratio is larger than one on next states whose probability RL increased and smaller than one
where RL suppressed probability.  It removes the donor's absolute base policy.

\subsection{Exact geometric tilt}

For transfer strength $\lambda>0$, define the conditional normalizer
\begin{equation}
Z_\lambda(z)
:=
\int
R_S(y\mid z)^\lambda
K_{\Lbase}(\dd y\mid z).
\label{eq:kernel-normalizer}
\end{equation}
For $\lambda=0$, define $R_S^0\equiv1$, so
$Z_0=1$ and $Q_0=K_{\Lbase}$ even on sets where $R_S=0$.

\begin{proposition}[KL-regularized donor-shift target]
\label{prop:kernel-tilt}
For $\lambda>0$, if $0<Z_\lambda(z)<\infty$, the kernel
\begin{equation}
\frac{
  \dd Q_\lambda(\cdot\mid z)
}{
  \dd K_{\Lbase}(\cdot\mid z)
}(y)
=
\frac{R_S(y\mid z)^\lambda}{Z_\lambda(z)}
\label{eq:kernel-tilt}
\end{equation}
is the unique maximizer of
\begin{equation}
\argmaxop_K
\left\{
\lambda\E_{Y\sim K}[r_S(Y\mid z)]
-
\KL(K\Vert K_{\Lbase})
\right\}.
\label{eq:kernel-variational}
\end{equation}
Equivalently, for every $K\ll K_{\Lbase}$ for which the two sides are
well-defined in the extended-real sense,
\begin{equation}
\lambda\E_K[r_S]
-\KL(K\Vert K_{\Lbase})
=
\log Z_\lambda
-\KL(K\Vert Q_\lambda).
\label{eq:kernel-kl-identity}
\end{equation}
If $K\not\ll Q_\lambda$, the right side and the variational objective are both interpreted as
$-\infty$ rather than as an indeterminate difference of infinities.
\end{proposition}

\begin{proof}
From \eqref{eq:kernel-tilt},
\[
\log\frac{\dd Q_\lambda}{\dd K_{\Lbase}}
=
\lambda r_S-\log Z_\lambda .
\]
Substitute this identity into
$\KL(K\Vert Q_\lambda)
=\E_K[\log(\dd K/\dd Q_\lambda)]$
and rearrange to obtain \eqref{eq:kernel-kl-identity}.  Nonnegativity of KL gives the maximizer.
\end{proof}

For densities with respect to a common base measure,
\begin{equation}
q_\lambda(y\mid z)
=
\frac{
  k_{\Lbase}(y\mid z)
  \left[
    k_{\Srl}(y\mid z)/k_{\Sbase}(y\mid z)
  \right]^\lambda
}{
  Z_\lambda(z)
}.
\label{eq:density-tilt}
\end{equation}
This is a \emph{geometric} or product-of-experts tilt.  It must not be confused with the
arithmetic probability mixture used by P-OPD,
\begin{equation}
(1-\alpha)K_{\mathrm{old}}+\alpha K_T.
\label{eq:not-mixture}
\end{equation}
The geometric target has no component-responsibility sigmoid.

\subsection{Implicit and non-Gaussian kernels}

The general construction is useful beyond Gaussian transitions whenever the ratio can be
estimated:
\begin{enumerate}[leftmargin=1.6em]
\item If all log densities are evaluable, optimize
      \begin{equation}
      \KL(K_{\Lactor}\Vert K_{\Lbase})
      -\lambda\E_{K_{\Lactor}}[r_S].
      \label{eq:implicit-kernel-loss}
      \end{equation}
      The unknown $\log Z_\lambda$ is independent of the trainable parameters.
\item If only donor transition samples are available, train a conditional discriminator with
      equal class priors.  Its Bayes optimum obeys
      \begin{equation}
      D^*(z,y)
      =
      \frac{k_{\Srl}(y\mid z)}
           {k_{\Srl}(y\mid z)+k_{\Sbase}(y\mid z)},
      \qquad
      \operatorname{logit}D^*(z,y)=r_S(y\mid z).
      \label{eq:ratio-classifier}
      \end{equation}
      Both classes must use the same query-state distribution; otherwise the classifier also
      learns a state-occupancy ratio.
\item One may sample approximately from $Q_\lambda$ by weighting
      $K_{\Lbase}$ transitions with $R_S^\lambda$, followed by resampling or SMC.  This is exact
      in the population limit but can have severe effective-sample-size collapse.
\item MMD, OT, score matching, Stein losses, or conditional flow matching can project implicit
      kernels into a trainable model.  Without an exact ratio identity they are surrogates, not
      exact realizations of \eqref{eq:kernel-tilt}.
\end{enumerate}

The kernel formulation therefore separates two statements:
\begin{itemize}[leftmargin=1.6em]
\item \emph{Flow matching does not require Gaussian probability paths.}
\item \emph{A chosen stochastic sampler may use Gaussian increments because they make the
      policy-shift ratio and target loss tractable.}
\end{itemize}

%% ============================================================
\section{Continuous-SDE formulation}
\label{sec:sde}
%% ============================================================

\subsection{Shared-diffusion processes and Girsanov ratio}

Use an increasing-time coordinate $\tau\in[0,T]$ and consider
\begin{equation}
\dd X_\tau
=
b_j(\tau,X_\tau,c)\dd\tau
+\sigma_\tau\dd W_\tau,
\qquad
a_\tau:=\sigma_\tau\sigma_\tau^\top\succ0.
\label{eq:continuous-sde}
\end{equation}
All donor and recipient processes share the diffusion coefficient and initial law; only their
drifts differ.  We assume the required Girsanov exponentials are true martingales for every pair
whose path measures are compared.  Define
\begin{equation}
\dbS:=b_{\Srl}-b_{\Sbase},
\qquad
u_\tau:=\sigma_\tau^\top a_\tau^{-1}\dbS.
\label{eq:sde-drift-shift}
\end{equation}

\begin{proposition}[Donor path ratio]
\label{prop:girsanov-ratio}
Under a Novikov or equivalent martingale condition, Girsanov's theorem gives
\begin{align}
\log\frac{\dd P_{\Srl}}{\dd P_{\Sbase}}
&=
\int_0^T u_\tau^\top\dd W_\tau^{\Sbase}
-\frac12\int_0^T\norm{u_\tau}^2\dd\tau
\label{eq:girsanov-brownian}\\
&=
\int_0^T
\dbS^\top a_\tau^{-1}
\left[
\dd X_\tau
-\frac{b_{\Srl}+b_{\Sbase}}{2}\dd\tau
\right].
\label{eq:girsanov-canonical}
\end{align}
All fields are evaluated along the same path $X_{0:T}$.
\end{proposition}

The corresponding path KL between two shared-diffusion processes is
\begin{equation}
\KL(P_b\Vert P_{\Lbase})
=
\frac12
\E_{P_b}
\int_0^T
\norm{b-b_{\Lbase}}_{a_\tau^{-1}}^2
\dd\tau .
\label{eq:girsanov-kl}
\end{equation}
This is an exact stochastic-process identity, not a Gaussian one-step
approximation~\cite{gradualft}.

\subsection{Local zero-discount control target}

At a frozen $(\tau,x,c)$, consider using a candidate drift $b$ over one infinitesimal interval.
The expected donor log-ratio rate under that candidate is
\begin{equation}
\mathfrak r_S(b)
=
\dbS^\top a^{-1}(b-b_{\Sbase})
-\frac12\norm{\dbS}_{a^{-1}}^2.
\label{eq:local-reward-rate}
\end{equation}
The recipient drift-KL rate relative to $\Lbase$ is
\begin{equation}
\mathfrak c_{\Lbase}(b)
=
\frac12\norm{b-b_{\Lbase}}_{a^{-1}}^2.
\label{eq:local-kl-rate}
\end{equation}

\begin{proposition}[Local SDE policy-shift graft]
\label{prop:local-sde-graft}
The frozen-state objective
\begin{equation}
\mathcal H_\lambda(b)
=
\lambda\mathfrak r_S(b)-\mathfrak c_{\Lbase}(b)
\label{eq:local-hamiltonian}
\end{equation}
has the unique optimum
\begin{equation}
\boxed{
b^*
=
b_{\Lbase}
+\lambda(b_{\Srl}-b_{\Sbase}).
}
\label{eq:local-drift-target}
\end{equation}
\end{proposition}

\begin{proof}
Terms independent of $b$ can be discarded.  Completing the square gives
\[
\mathcal H_\lambda(b)
=
\mathrm{constant}
-\frac12
\norm{
b-[b_{\Lbase}+\lambda\dbS]
}_{a^{-1}}^2 .
\]
Positive definiteness of $a$ gives the unique optimum.
\end{proof}

Equation~\eqref{eq:local-drift-target} is the continuous-SDE counterpart of adding a small-model
logit delta to a large language model.  It is exact for the frozen-state, immediate-reward
problem.  Using recipient on-policy states but detaching them during optimization produces a
semi-gradient version of this objective.

\subsection{Why the globally tilted path needs a value correction}

For this subsection, additionally assume
\begin{equation}
P_{\Lbase}\ll P_{\Sbase},
\qquad
0<\mathcal Z_\lambda
:=
\E_{P_{\Lbase}}
\left[
\left(
\frac{\dd P_{\Srl}}{\dd P_{\Sbase}}
\right)^\lambda
\right]
<\infty .
\label{eq:path-tilt-assumptions}
\end{equation}
The ideal path-space target is
\begin{equation}
\frac{\dd\mathbb Q_\lambda}{\dd P_{\Lbase}}
=
\frac1{\mathcal Z_\lambda}
\left(
\frac{\dd P_{\Srl}}{\dd P_{\Sbase}}
\right)^\lambda.
\label{eq:global-path-tilt}
\end{equation}
Define
\begin{equation}
\bar b=b_{\Lbase}+\lambda\dbS
\label{eq:myopic-drift}
\end{equation}
and the state-dependent potential
\begin{equation}
U_\lambda
=
\lambda\dbS^\top a^{-1}(b_{\Lbase}-b_{\Sbase})
+\frac{\lambda(\lambda-1)}2
\norm{\dbS}_{a^{-1}}^2.
\label{eq:path-potential}
\end{equation}
Completing the stochastic exponential shows that
\eqref{eq:global-path-tilt} is a Feynman--Kac weighting of the process with drift $\bar b$.
Let
\begin{equation}
H(\tau,x)
=
\E_{\bar b}^{\tau,x}
\left[
\exp\left(
\int_\tau^T U_\lambda(s,X_s)\dd s
\right)
\right].
\label{eq:desirability}
\end{equation}
Under the usual regularity conditions, $H$ solves
\begin{equation}
\partial_\tau H
+\bar b\cdot\nabla H
+\frac12a:\nabla^2H
+U_\lambda H
=0,
\qquad
H(T,x)=1.
\label{eq:desirability-pde}
\end{equation}
The globally tilted process has Doob-transformed drift
\begin{equation}
\boxed{
b_{\mathbb Q_\lambda}
=
b_{\Lbase}
+\lambda(b_{\Srl}-b_{\Sbase})
+a\nabla\log H.
}
\label{eq:global-drift}
\end{equation}
If the shared initial law is $\rho_0$, the unconditional global tilt also changes it:
\begin{equation}
\mathbb Q_\lambda(X_0\in\dd x)
=
\frac{H(0,x)\rho_0(\dd x)}{\mathcal Z_\lambda},
\qquad
\mathcal Z_\lambda
=
\int H(0,x)\rho_0(\dd x).
\label{eq:tilted-initial-law}
\end{equation}
Therefore the drift \eqref{eq:global-drift} started from the original $\rho_0$ does not by itself
realize the unconditional target \eqref{eq:global-path-tilt}.  If the generative base-noise law
must remain fixed, define the tilt conditionally on $X_0=x$ and normalize each conditional path
law by $H(0,x)$ before averaging over the original $\rho_0$.

The last term is a continuation-value correction.  It vanishes only in special cases, for example
when $U_\lambda$ is spatially constant.  Therefore simple drift arithmetic should be described as
the local, zero-discount Direct-OPD target, not as the exact globally normalized path tilt.  This
is the continuous analogue of Direct-OPD's practical immediate-token surrogate versus its
idealized sequence-level objective.

%% ============================================================
\section{Gaussian, Euler-SDE, and DPM-SDE specializations}
\label{sec:dpm-sde}
%% ============================================================

\subsection{Generic shared-covariance Gaussian transition}

Condition on a recipient query state $z$.  Suppose the donor pair shares covariance $\Sigma_S$
and the recipient pair shares covariance $\Sigma_L$:
\begin{align}
K_{\Sbase}&=\N(\mu_{\Sbase},\Sigma_S),
&
K_{\Srl}&=\N(\mu_{\Srl},\Sigma_S),
\nonumber\\
K_{\Lbase}&=\N(\mu_{\Lbase},\Sigma_L),
&
K_{\Lactor}&=\N(\mu_\theta,\Sigma_L).
\label{eq:two-covariance-gaussians}
\end{align}
Define
\begin{equation}
\dmuS:=\mu_{\Srl}-\mu_{\Sbase},
\qquad
\bar\mu_S:=\frac{\mu_{\Srl}+\mu_{\Sbase}}2.
\label{eq:mean-shift}
\end{equation}
The donor log ratio is linear in $y$:
\begin{equation}
r_S(y)
=
\dmuS^\top\Sigma_S^{-1}(y-\bar\mu_S).
\label{eq:gaussian-log-ratio}
\end{equation}

\begin{proposition}[Closed-form Gaussian geometric tilt]
\label{prop:gaussian-tilt}
The kernel target \eqref{eq:kernel-tilt} is
\begin{equation}
Q_\lambda
=
\N(\mu_\lambda^*,\Sigma_L),
\qquad
\boxed{
\mu_\lambda^*
=
\mu_{\Lbase}
+\lambda\Sigma_L\Sigma_S^{-1}\dmuS.
}
\label{eq:gaussian-target-mean}
\end{equation}
Consequently,
\begin{equation}
\boxed{
\KL(K_{\Lactor}\Vert Q_\lambda)
=
\frac12
\norm{
\mu_\theta-\mu_{\Lbase}
-\lambda\Sigma_L\Sigma_S^{-1}\dmuS
}_{\Sigma_L^{-1}}^2.
}
\label{eq:gaussian-training-kl}
\end{equation}
\end{proposition}

\begin{proof}
Insert \eqref{eq:gaussian-log-ratio} into
$\log K_{\Lbase}(y)+\lambda r_S(y)$.  Its quadratic term remains
$-\tfrac12 y^\top\Sigma_L^{-1}y$ and its linear term is
\[
y^\top
\left[
\Sigma_L^{-1}\mu_{\Lbase}
+\lambda\Sigma_S^{-1}\dmuS
\right].
\]
Completing the square gives \eqref{eq:gaussian-target-mean}.  The equal-covariance Gaussian KL
from $K_{\Lactor}$ to $Q_\lambda$ gives \eqref{eq:gaussian-training-kl}.
\end{proof}

If all four covariances are the same,
\begin{equation}
\boxed{
\mu_\lambda^*
=
\mu_{\Lbase}
+\lambda(\mu_{\Srl}-\mu_{\Sbase}).
}
\label{eq:equal-covariance-target}
\end{equation}
Thus mean-delta arithmetic is not an arbitrary heuristic in this regime; it is the exact natural
parameter update selected by the general kernel objective.

\begin{remark}[Unequal donor covariances]
If $\Sbase$ and $\Srl$ do not share covariance, the donor reward is quadratic and the target
precision becomes
\begin{equation}
\Lambda_*
=
\Lambda_{\Lbase}
+\lambda(\Lambda_{\Srl}-\Lambda_{\Sbase}),
\label{eq:unequal-precision}
\end{equation}
which is normalizable only if $\Lambda_*\succ0$.  Plain mean or velocity subtraction is then not
the exact target.
\end{remark}

\subsection{Euler--Maruyama SDE}

For a step of length $h>0$, let
\begin{equation}
K_j^h(\cdot\mid x)
=
\N(x+h\,b_j,\;hA),
\qquad
A\succ0.
\label{eq:euler-sde-kernel}
\end{equation}
All drifts are evaluated at the same $(t,x,c)$.  Proposition~\ref{prop:gaussian-tilt} gives
\begin{equation}
Q_\lambda^h
=
\N\!\left(
x+h[b_{\Lbase}+\lambda(b_{\Srl}-b_{\Sbase})],
hA
\right).
\label{eq:euler-target}
\end{equation}
Writing
\begin{equation}
b_*:=b_{\Lbase}+\lambda(b_{\Srl}-b_{\Sbase}),
\end{equation}
the exact event-sum conditional KL is
\begin{equation}
\ell_{\mathrm{Euler,sum}}
=
\frac{h}{2}
\norm{b_{\Lactor}-b_*}_{A^{-1}}^2.
\label{eq:euler-sde-kl}
\end{equation}
For isotropic $A=g_t^2I$ and event dimension $D$,
\begin{align}
\ell_{\mathrm{Euler,sum}}
&=
\frac{h}{2g_t^2}
\sum_{d=1}^D
(b_{\Lactor,d}-b_{*,d})^2,
\label{eq:euler-sum}\\
\ell_{\mathrm{Euler,mean}}
&=
\frac1D\ell_{\mathrm{Euler,sum}}.
\label{eq:euler-mean}
\end{align}
The sum is the joint transition KL; the mean is its dimension-normalized version.  They have the
same isolated minimizer but different strength relative to other losses and trust budgets.

\subsection{First-order SDE-DPM-Solver++}

SDE-DPM-Solver++ is a stochastic solver for a \emph{chosen Gaussian diffusion path}; it is not a
generic property of arbitrary flow matching.  Assume
\begin{equation}
X_t=\alpha_tX_0+\sigma_t\varepsilon,
\qquad
\varepsilon\sim\N(0,I),
\qquad
\ell_t:=\log\frac{\alpha_t}{\sigma_t}.
\label{eq:dpm-path}
\end{equation}
For a reverse step from noisier $s$ to cleaner $t$, set
\begin{equation}
h:=\ell_t-\ell_s>0.
\label{eq:dpm-h}
\end{equation}
The first-order SDE-DPM-Solver++ update~\cite{dpmsolverpp} is
\begin{equation}
X_t
=
\underbrace{\frac{\sigma_t}{\sigma_s}e^{-h}}_{A_{t,s}}X_s
+
\underbrace{\alpha_t(1-e^{-2h})}_{B_{t,s}}
\widehat x_{0,j}(X_s,s)
+
\underbrace{\sigma_t\sqrt{1-e^{-2h}}}_{C_{t,s}}\xi,
\qquad
\xi\sim\N(0,I).
\label{eq:dpm-first-order}
\end{equation}
Conditioned on $X_s=x_s$,
\begin{equation}
K_j(\cdot\mid x_s)
=
\N(\mu_j,\Sigma_{t,s}),
\quad
\mu_j=A_{t,s}x_s+B_{t,s}\widehat x_{0,j},
\quad
\Sigma_{t,s}=C_{t,s}^2I.
\label{eq:dpm-first-kernel}
\end{equation}
The stochastic covariance is model independent and the mean is affine in the model's data
prediction.  Therefore the exact graft is
\begin{align}
\mu_*
&=
\mu_{\Lbase}
+\lambda(\mu_{\Srl}-\mu_{\Sbase})
\nonumber\\
&=
A_{t,s}x_s
+B_{t,s}
\left[
\widehat x_{0,\Lbase}
+\lambda(
\widehat x_{0,\Srl}
-\widehat x_{0,\Sbase}
)
\right],
\label{eq:dpm-grafted-mean}
\end{align}
or equivalently
\begin{equation}
\boxed{
\widehat x_{0,*}
=
\widehat x_{0,\Lbase}
+\lambda(
\widehat x_{0,\Srl}
-\widehat x_{0,\Sbase}
).
}
\label{eq:dpm-x0-target}
\end{equation}
The exact solver-kernel loss is
\begin{equation}
\KL(K_{\Lactor}\Vert Q_\lambda)
=
\frac{B_{t,s}^2}{2C_{t,s}^2}
\norm{\widehat x_{0,\Lactor}-\widehat x_{0,*}}^2.
\label{eq:dpm-first-kl}
\end{equation}
This equality is exact for the selected first-order numerical transition kernel.  It does not
claim that the numerical solver exactly integrates the underlying continuous SDE.

\subsection{Multistep SDE-DPM-Solver++}

A multistep solver is Markov only after its history is included in the state.  For
SDE-DPM-Solver++(2M), let $r$ be the previous history point and define
\begin{equation}
r_0:=\frac{\ell_s-\ell_r}{h}>0,
\quad
D_{0,j}:=\widehat x_{0,j}(x_s,s),
\quad
D_{1,j}:=
\frac{
\widehat x_{0,j}(x_s,s)
-\widehat x_{0,j}(x_r,r)
}{r_0}.
\label{eq:dpm-history-terms}
\end{equation}
Equation (20) of DPM-Solver++ can be written
\begin{equation}
\mu_j
=
A_{t,s}x_s
+B_{t,s}D_{0,j}
+\frac{B_{t,s}}2D_{1,j},
\qquad
\Sigma_{t,s}=C_{t,s}^2I.
\label{eq:dpm-2m-mean}
\end{equation}
The exact target conditioned on the common augmented history is
\begin{equation}
\boxed{
\mu_*
=
\mu_{\Lbase}
+\lambda(\mu_{\Srl}-\mu_{\Sbase}).
}
\label{eq:dpm-2m-target}
\end{equation}
Every current and historical model-output term contributes to the mean difference.  Grafting only
the current velocity or current $\widehat x_0$ omits the multistep derivative correction and is
not the exact conditional target.

More generally, if
\begin{equation}
\mu_j(\mathcal H_n)
=
A_nx_n+\sum_{q=0}^{m-1}B_{nq}M_{j,nq}(\mathcal H_n),
\label{eq:generic-multistep}
\end{equation}
then compute the \emph{whole solver mean} for each model on the same history and graft
$\mu_{\Lbase}+\lambda(\mu_{\Srl}-\mu_{\Sbase})$.  Marginalizing random histories generally
produces a non-Gaussian mixture kernel.

\subsection{When a flow velocity can use DPM-SDE}

For a known affine Gaussian conditional path
\begin{equation}
X_t=\alpha_tX_0+\sigma_t\varepsilon,
\qquad
\dot X_t=\dot\alpha_tX_0+\dot\sigma_t\varepsilon,
\label{eq:affine-path}
\end{equation}
define
\begin{equation}
\mathcal D_t
:=
\alpha_t\dot\sigma_t-\sigma_t\dot\alpha_t
\neq0.
\label{eq:path-determinant}
\end{equation}
A flow-velocity prediction $v_j(t,x)$ converts to
\begin{align}
\widehat x_{0,j}
&=
\frac{\dot\sigma_t x-\sigma_tv_j(t,x)}{\mathcal D_t},
\label{eq:velocity-to-x0}\\
\widehat\varepsilon_j
&=
\frac{-\dot\alpha_t x+\alpha_tv_j(t,x)}{\mathcal D_t}.
\label{eq:velocity-to-epsilon}
\end{align}
For normalized VP schedules satisfying
$\alpha_t^2+\sigma_t^2=1$, standard diffusion $v$-prediction defines
\begin{equation}
v_{\mathrm{pred}}
=
\alpha_t\varepsilon-\sigma_tX_0,
\qquad
\widehat x_0
=
\alpha_tx-\sigma_tv_{\mathrm{pred}}.
\label{eq:standard-v-pred}
\end{equation}
The physical path derivative and the standard network output differ by
\begin{equation}
v_{\mathrm{physical}}
=
\mathcal D_t\,v_{\mathrm{pred}}.
\label{eq:physical-v-scale}
\end{equation}
Thus a raw $v_{\mathrm{pred}}$ must not be substituted for the physical flow velocity without this
schedule-dependent conversion.

DPM-SDE is applicable only when the flow model's parameterization belongs to a compatible
Gaussian diffusion schedule with monotone half-logSNR and a valid reverse SDE.  An arbitrary
non-Gaussian flow-matching path, nonlinear interpolation, or bare learned ODE field does not
canonically define $\widehat x_0$, a score, a reverse diffusion, or a DPM-SDE transition.

\begin{tcolorbox}[colback=red!4,colframe=red!45!black,title=DPM-SDE exactness boundary]
The final zero-variance DPM step is deterministic and has no transition-density ratio.  Exact
kernel grafting requires positive model-independent covariance, a common augmented solver history,
and model outputs converted into the same prediction parameterization.
\end{tcolorbox}

%% ============================================================
\section{Deterministic ODE formulations}
\label{sec:ode}
%% ============================================================

\subsection{Why transition ratios fail}

For an ODE Euler step,
\begin{equation}
K_j(\dd y\mid x)
=
\delta_{x+h\,v_j(t,x,c)}(\dd y).
\label{eq:ode-dirac}
\end{equation}
If the donor maps differ, $K_{\Srl}$ and $K_{\Sbase}$ are mutually singular, so
$\dd K_{\Srl}/\dd K_{\Sbase}$ is undefined.  Girsanov also fails because the diffusion covariance
is zero.  No deterministic transition KL or policy ratio can justify the target directly.

\subsection{Selected small-noise control limit}

Introduce a reference small-noise family
\begin{equation}
K_j^{\varepsilon,h}
=
\N(x+h\,v_j,\;\varepsilon hA),
\qquad
A\succ0.
\label{eq:small-noise-family}
\end{equation}
For every $\varepsilon>0$, the exact geometric target has mean
\begin{equation}
x+h
\left[
v_{\Lbase}
+\lambda(v_{\Srl}-v_{\Sbase})
\right].
\label{eq:small-noise-target}
\end{equation}
Define
\begin{equation}
v_*:=v_{\Lbase}+\lambda(v_{\Srl}-v_{\Sbase}).
\label{eq:ode-field-target}
\end{equation}
Then
\begin{equation}
\frac{\varepsilon}{h}
\KL(K_{\Lactor}^{\varepsilon,h}\Vert Q_\lambda^{\varepsilon,h})
=
\frac12
\norm{v_{\Lactor}-v_*}_{A^{-1}}^2.
\label{eq:ode-control-limit}
\end{equation}
The ratio itself does not exist at $\varepsilon=0$, but the Gaussian family selects a finite
deterministic target and its rescaled KL is exactly a quadratic control metric.

The practical recipient-on-policy ODE loss is therefore
\begin{equation}
\boxed{
\mathcal L_{\mathrm{ODE}}
=
\E_{\substack{
C,t\\
X_t\sim\rho_{\Lbeh,t}
}}
\left[
w(t)
\norm{
[v_{\Lactor}-v_{\Lbase}]
-\lambda[v_{\Srl}-v_{\Sbase}]
}_{G_t}^2
\right].
}
\label{eq:ode-main-loss}
\end{equation}
Without a declared small-noise SDE, $G_t$ and $w(t)$ are modeling choices.  Common options are
identity velocity geometry, a scheduler-derived control metric, or a per-dimension RMS
normalization.

\subsection{Lagrangian flow-map delta}

The Eulerian field loss queries the 9B donor on 32B recipient states.  If those states are too far
outside the donor occupancy, a Lagrangian alternative evaluates every model on its own trajectory.
Let $\Phi_{j,t}(Z,c)$ be flow maps from common base noise $Z$ and define
\begin{equation}
\Psi_t^\lambda(Z,c)
=
\Phi_{\Lbase,t}(Z,c)
+\lambda
\left[
\Phi_{\Srl,t}(Z,c)
-\Phi_{\Sbase,t}(Z,c)
\right].
\label{eq:flow-map-delta}
\end{equation}
Its path derivative is
\begin{align}
\dot\Psi_t^\lambda
&=
v_{\Lbase}(t,\Phi_{\Lbase,t},c)
\nonumber\\
&\quad+
\lambda
\left[
v_{\Srl}(t,\Phi_{\Srl,t},c)
-v_{\Sbase}(t,\Phi_{\Sbase,t},c)
\right].
\label{eq:flow-map-velocity}
\end{align}
Conditional flow matching uses
\begin{equation}
\mathcal L_{\mathrm{map\text{-}CFM}}
=
\E_{C,Z,t}
\norm{
v_\theta(t,\Psi_t^\lambda,C)
-\dot\Psi_t^\lambda
}^2.
\label{eq:flow-map-cfm}
\end{equation}
Its population optimum is the conditional mean
\begin{equation}
u_t(x,c)
=
\E[
\dot\Psi_t^\lambda(Z,c)
\mid
\Psi_t^\lambda(Z,c)=x
].
\label{eq:map-marginal-field}
\end{equation}
Under standard regularity, this field transports the chosen pseudo-path
marginals~\cite{flowmatching}.
It is exact for the explicitly declared common-noise coupling, not a zero-noise policy-ratio
identity.

\subsection{Marginal density-ratio target}

ODE time marginals may possess smooth densities even though their transition kernels are Dirac.
Assume, at each $(t,c)$,
\begin{equation}
\rho_{\Srl,t}(\cdot\mid c)\ll\rho_{\Sbase,t}(\cdot\mid c),
\qquad
\rho_{\Lbase,t}(\cdot\mid c)\ll\rho_{\Sbase,t}(\cdot\mid c),
\label{eq:marginal-support}
\end{equation}
and define the finite positive normalizer
\begin{equation}
Z_t(c)
=
\int
\rho_{\Lbase,t}(x\mid c)
\left[
\frac{\rho_{\Srl,t}(x\mid c)}
     {\rho_{\Sbase,t}(x\mid c)}
\right]^\lambda
\dd x
\in(0,\infty).
\label{eq:marginal-normalizer}
\end{equation}
The static density target is
\begin{equation}
q_t(x\mid c)
=
\frac1{Z_t(c)}
\rho_{\Lbase,t}(x\mid c)
\left[
\frac{\rho_{\Srl,t}(x\mid c)}
     {\rho_{\Sbase,t}(x\mid c)}
\right]^\lambda.
\label{eq:marginal-ratio-target}
\end{equation}
On the region where all three component densities are positive, its static score identity is
\begin{equation}
\nabla\log q_t
=
\nabla\log\rho_{\Lbase,t}
+\lambda
\left[
\nabla\log\rho_{\Srl,t}
-\nabla\log\rho_{\Sbase,t}
\right].
\label{eq:marginal-score-arithmetic}
\end{equation}
However, products and ratios of densities following their respective continuity or
Fokker--Planck dynamics do not generally follow the dynamics induced by score arithmetic.
Consequently, score or velocity arithmetic does not automatically generate
the complete path \eqref{eq:marginal-ratio-target}.  A strict endpoint version requires an
endpoint density-ratio estimator, importance resampling, DMD with an online fake score, or a
dynamic corrector.  This is a separate and more expensive distribution-level method.

%% ============================================================
\section{Main training objective and algorithm}
\label{sec:algorithm}
%% ============================================================

\subsection{Recommended transition-mean implementation}

The safest implementation uses the \emph{actual scheduler or solver transition mean} rather than
assuming a particular Euler coefficient.  On a recipient behavior state and common solver
history, compute
\begin{equation}
\dmuS
=
\mu_{\Srl}-\mu_{\Sbase}.
\label{eq:actual-mean-shift}
\end{equation}
When donor and recipient covariances differ but each pair is internally matched, define
\begin{equation}
d_\mu
:=
\Sigma_L\Sigma_S^{-1}\dmuS.
\label{eq:preconditioned-shift}
\end{equation}
The unit-strength recipient trust cost is
\begin{equation}
K_1
:=
\frac12
d_\mu^\top\Sigma_L^{-1}d_\mu.
\label{eq:unit-trust-kl}
\end{equation}
Define the corresponding per-dimension quantities
\begin{equation}
\bar K_1:=K_1/D,
\qquad
\bar\kappa:=\kappa/D.
\label{eq:normalized-trust}
\end{equation}
Given nominal strength $\lambda_0$ and either a conditional joint-KL budget $\kappa$ or the
equivalent per-dimension budget $\bar\kappa$, set
\begin{equation}
\lambda_{\mathrm{eff}}
=
\begin{cases}
\lambda_0, & K_1=0,\\[0.3em]
\min\!\left\{
\lambda_0,\sqrt{\kappa/K_1}
\right\}
=
\min\!\left\{
\lambda_0,\sqrt{\bar\kappa/\bar K_1}
\right\},
& K_1>0.
\end{cases}
\label{eq:lambda-clipping}
\end{equation}
This guarantees
\begin{equation}
\KL(Q_{\lambda_{\mathrm{eff}}}\Vert K_{\Lbase})
\le\kappa .
\label{eq:target-trust-guarantee}
\end{equation}
The target mean is
\begin{equation}
\mu_{\mathrm{target}}
=
\mu_{\Lbase}
+\lambda_{\mathrm{eff}}d_\mu.
\label{eq:main-target-mean}
\end{equation}

The dimension-normalized training loss is
\begin{equation}
\boxed{
\mathcal L_{\mathrm{Flow\text{-}Direct\text{-}OPD}}
=
\E_{z\sim\nu_{\Lbeh}}
\left[
\frac{1}{2D}
\norm{
\mu_\theta(z)
-\sg(\mu_{\mathrm{target}}(z))
}_{\Sigma_L^{-1}}^2
\right].
}
\label{eq:main-training-loss}
\end{equation}
Removing $1/D$ recovers the exact joint conditional KL.  The displayed optimization loss uses the
per-dimension convention; \eqref{eq:target-trust-guarantee} is the equivalent joint guarantee.
Equation~\eqref{eq:normalized-trust} makes the conversion explicit.

If one shared isotropic Euler-SDE wrapper is used, this simplifies to
\begin{align}
v_{\mathrm{target}}
&=
v_{\Lbase}
+\lambda_{\mathrm{eff}}
(v_{\Srl}-v_{\Sbase}),
\label{eq:main-velocity-target}\\
\mathcal L
&=
\E_{\nu_{\Lbeh}}
\left[
\frac{h_t}{2Dg_t^2}
\norm{
v_{\Lactor}-\sg(v_{\mathrm{target}})
}^2
\right].
\label{eq:main-velocity-loss}
\end{align}
For the ODE arm, use \eqref{eq:ode-main-loss} instead of claiming a transition KL.

\subsection{Recipient-on-policy optimization}

\begin{algorithm}[htbp]
\caption{Weak-to-Strong Flow Direct-OPD}
\label{alg:flow-direct-opd}
\begin{algorithmic}[1]
\Require frozen donor reference $\Sbase$; frozen donor RL checkpoint $\Srl$
\Require large recipient reference $\Lbase$; trainable recipient $\Lactor\leftarrow\Lbase$
\Require nominal strength $\lambda_0$; optional trust budget $\kappa$
\For{outer iteration $k=0,1,\ldots$}
  \State Freeze behavior checkpoint $\Lbeh\leftarrow\Lactor$
  \State Roll out $\Lbeh$ and cache recipient states, physical times, and solver histories
  \For{each sampled training state/history $z$}
    \State Evaluate $\mu_{\Sbase}(z)$ and $\mu_{\Srl}(z)$ without gradients
    \State Evaluate $\mu_{\Lbase}(z)$ without gradients
    \State Compute $\dmuS$, $d_\mu$, and $K_1$
    \State If $\kappa$ is absent set $\lambda_{\mathrm{eff}}\leftarrow\lambda_0$;
      otherwise use \eqref{eq:lambda-clipping}
    \State Set $\mu_{\mathrm{target}}\leftarrow
      \mu_{\Lbase}+\lambda_{\mathrm{eff}}d_\mu$
    \State Evaluate trainable $\mu_\theta(z)$ and accumulate
      \eqref{eq:main-training-loss}
  \EndFor
  \State Update $\theta$; log realized recipient-to-base displacement and refresh behavior
\EndFor
\end{algorithmic}
\end{algorithm}

The donor pair and recipient reference are detached.  The query distribution is the recipient's
own behavior occupancy, which is the flow analogue of Direct-OPD evaluating the weak checkpoint
ratio on strong-student prefixes.  Refreshing the behavior model limits stale-state error, but the
detached optimization remains a local semi-gradient rather than the full path gradient of
\eqref{eq:global-path-tilt}.

\subsection{Reliability masks and composition}

If the donor shift is unreliable on some recipient states, introduce a detached mask
$m(z)\in[0,1]$:
\begin{equation}
\mu_{\mathrm{target}}
=
\mu_{\Lbase}
+m(z)\lambda_{\mathrm{eff}}d_\mu.
\label{eq:masked-target}
\end{equation}
Possible masks use cross-seed or cross-checkpoint donor-delta agreement, donor-state OOD tests, or
explicit domain routing.  A mask based only on $\norm{\dmuS}$ is not a correctness test: a large
shift can encode a reward shortcut.

Several independently learned shifts can be composed:
\begin{equation}
\mu_{\mathrm{target}}
=
\mu_{\Lbase}
+\sum_{m=1}^M
\lambda_m d_\mu^{(m)}.
\label{eq:composed-shifts}
\end{equation}
This is exact for a product of conditional donor ratios when the normalizer exists.  In practice,
joint trust projection is needed because individually safe shifts can add constructively or
conflict.

%% ============================================================
\section{Alternative objectives and their exactness}
\label{sec:alternatives}
%% ============================================================

\begin{center}
\begin{tabular}{@{}
  >{\raggedright\arraybackslash}p{0.05\textwidth}
  >{\raggedright\arraybackslash}p{0.19\textwidth}
  >{\raggedright\arraybackslash}p{0.27\textwidth}
  >{\raggedright\arraybackslash}p{0.37\textwidth}
@{}}
\toprule
arm & construction & exact meaning & principal boundary\\
\midrule
A & inference-time velocity graft
& a declared ODE field; exact conditional Gaussian mean under shared SDE covariance
& requires all three frozen models at inference and can move donor queries off-support\\
\addlinespace
B & recipient-on-policy residual field distillation
& weighted $L^2$ projection of arm A under the behavior occupancy
& finite capacity and stale occupancy; not an endpoint distribution objective\\
\addlinespace
C & kernel or SDE transition tilt
& exact local KL-regularized target \eqref{eq:kernel-tilt}
& needs density ratio or shared stochastic kernel; global path target has Doob correction\\
\addlinespace
D & endpoint ratio tilt
& exact static distribution target \eqref{eq:marginal-ratio-target}
& expensive ratio/score estimation; product path is not automatically dynamically consistent\\
\addlinespace
E & common-noise flow-map CFM
& exact CFM target for the explicit pseudo-path \eqref{eq:flow-map-delta}
& coupling- and coordinate-dependent; additive maps can fold or leave the valid region\\
\bottomrule
\end{tabular}
\end{center}

\subsection{Zero-training inference oracle}

Before any recipient training, integrate
\begin{equation}
\dot X_t
=
v_{\Lbase}(t,X_t,c)
+\lambda
\left[
v_{\Srl}(t,X_t,c)
-v_{\Sbase}(t,X_t,c)
\right].
\label{eq:inference-graft}
\end{equation}
This needs one 32B and two 9B forwards per solver query, but directly tests whether the field
direction transfers.  If a correctly implemented closed-loop graft provides no
reward--quality Pareto improvement, training arm B toward the same field target is unlikely to
rescue the hypothesis.

\subsection{Endpoint importance-CFM}

Under the support and normalizability assumptions
\eqref{eq:marginal-support}--\eqref{eq:marginal-normalizer} at $t=1$, estimate
\begin{equation}
r_1(y,c)
=
\log\rho_{\Srl,1}(y\mid c)
-\log\rho_{\Sbase,1}(y\mid c).
\label{eq:endpoint-reward}
\end{equation}
Generate $Y_i\sim\rho_{\Lbase,1}$ and weight
\begin{equation}
w_i=\exp(\lambda r_1(Y_i,c)).
\label{eq:endpoint-weights}
\end{equation}
The population weight normalizer is
$Z_1(c)=\E_{Y\sim\rho_{\Lbase,1}}[w(Y,c)]\in(0,\infty)$.
Resampling by $w_i$ followed by ordinary endpoint CFM targets
\begin{equation}
q_1(y\mid c)
\propto
\rho_{\Lbase,1}(y\mid c)e^{\lambda r_1(y,c)}
\label{eq:endpoint-target}
\end{equation}
in the population limit.  Its feasibility is governed by
\begin{equation}
\mathrm{ESS}
=
\frac{(\sum_iw_i)^2}{\sum_iw_i^2}.
\label{eq:ess}
\end{equation}
Collapsed ESS is an estimation failure, not evidence that policy-shift transfer itself is false.
DMD~\cite{dmd} is an alternative when an online fake-score model is affordable.

\subsection{Comparison with direct teacher imitation and P-OPD}

Vanilla diffusion/flow distillation~\cite{diffusionopd} uses
$v_{\mathrm{target}}=v_{\Srl}$ and transfers the donor's absolute policy.  The canonical case in
which it agrees with Flow Direct-OPD is $\Sbase=\Lbase$ and $\lambda=1$; more generally equality
requires
$v_{\Lbase}+\lambda(v_{\Srl}-v_{\Sbase})=v_{\Srl}$
pointwise.

P-OPD uses an arithmetic mixture and a posterior responsibility.  Its exact local-gradient
interpretation requires a sampled stochastic transition and positive shared covariance; the gate
collapses in the deterministic limit.  Flow Direct-OPD uses a geometric ratio tilt.  Under shared
Gaussian covariance the product closes to one Gaussian and no responsibility gate appears.  The
two methods solve different proximal problems.

%% ============================================================
\section{FLUX.2 compatibility, diagnostics, and experiments}
\label{sec:practice}
%% ============================================================

\subsection{What the public model configs establish}

The public Hugging Face configs for
\href{https://huggingface.co/black-forest-labs/FLUX.2-klein-base-9B}{FLUX.2-klein-base-9B}
and
\href{https://huggingface.co/black-forest-labs/FLUX.2-dev}{FLUX.2-dev}
show encouraging structural compatibility:
\begin{center}
\begin{tabular}{@{}
  >{\raggedright\arraybackslash}p{0.17\textwidth}
  >{\raggedright\arraybackslash}p{0.35\textwidth}
  >{\raggedright\arraybackslash}p{0.37\textwidth}
@{}}
\toprule
component & observed relation & implication\\
\midrule
scheduler & same FlowMatchEulerDiscreteScheduler fields & physical sigma grids can be matched\\
transformer I/O & Flux2Transformer2DModel, 128 channels & output tensor geometry can match\\
RoPE axes & $(32,32,32,32)$ for both & latent token coordinate convention is compatible\\
VAE config & AutoencoderKLFlux2, 32 latent channels & same declared latent architecture\\
text encoder & Qwen3 vs.\ Mistral3 & embeddings are not interchangeable\\
joint attention dim & $7680$ vs.\ $15360$ & each model must use its own conditioning path\\
\bottomrule
\end{tabular}
\end{center}

Configuration equality is not numerical identity.  The two published VAE weight files have
different sizes (approximately $168.1$MB and $336.2$MB), consistent with a precision difference
but not proving equal fp32 values.  Train-time identity grafting requires direct state and
one-step parity tests.

\subsection{Fail-closed compatibility checks}

Before training:
\begin{enumerate}[leftmargin=1.6em]
\item Pin the exact $\Sbase$ revision from which the RL checkpoint was produced.  Disabling the
      RL adapter must recover that base numerically.
\item Compare VAE state dictionaries after fp32 conversion; encode the same images with posterior
      mode and compare packed latent coordinates, latent IDs, scaling, and decode parity.
\item Use one diffusers revision and compare complete sigma/timestep grids, dynamic shift, and
      solver coefficients at every intended resolution and NFE.
\item Instrument one step of each official pipeline and verify that converted outputs represent
      the same physical velocity or complete transition mean, including sign.
\item Feed the same raw prompt to each model's own text encoder.  Begin with guidance scale $1$;
      equal CFG numbers across the two pipelines need not define equal physical fields.
\item Compute donor subtraction in fp32 and deterministic eval mode.  Require its RMS to be well
      above repeated-forward and dtype noise before interpreting it as an RL signal.
\end{enumerate}

If latent coordinates differ, a state transport map $T$ is required and vector shifts must be
pushforwarded by its Jacobian.  Matching tensor shape alone is insufficient.

\subsection{Required scale and reliability logs}

At every sampled timestep, log both RMS and event-L2 scales:
\begin{itemize}[leftmargin=1.6em]
\item $\operatorname{RMS}(\dvS)$,
      $\operatorname{RMS}(v_{\Lbase})$, and their ratio;
\item one-step displacement $|h_t|\operatorname{RMS}(\lambda\dvS)$;
\item joint and per-dimension $K_1$, $\lambda_{\mathrm{eff}}$, and clipping rate;
\item open-loop and closed-loop endpoint latent displacement;
\item cross-seed, cross-checkpoint, or cross-scale cosine similarity of donor deltas;
\item recipient-state OOD diagnostics relative to $\Sbase$ and $\Srl$ trajectory states;
\item realized post-update
      $\operatorname{RMS}(v_{\Lactor}-v_{\Lbase})$
      versus the requested target.
\end{itemize}

For donor pairs $i,j$, let $G_k=A_k^{-1}$ for an SDE control metric, or use the explicitly
declared ODE metric.  A control-energy weighted cosine is
\begin{equation}
\operatorname{Cos}_{ij}
=
\frac{
\sum_k h_k
(\Delta v_{S,k}^{(i)})^\top G_k\Delta v_{S,k}^{(j)}
}{
\sqrt{\sum_k h_k\norm{\Delta v_{S,k}^{(i)}}_{G_k}^2}
\sqrt{\sum_k h_k\norm{\Delta v_{S,k}^{(j)}}_{G_k}^2}
}.
\label{eq:delta-cosine}
\end{equation}
This measures signal stability, not agreement with an unavailable 32B RL delta.

\subsection{Minimum informative experiment sequence}

\paragraph{Stage 0: compatibility.}
Use $32$--$64$ prompts across intended resolutions and timesteps.  Stop on latent, scheduler,
parameterization, conditioning, or checkpoint-provenance failure.

\paragraph{Stage 1: open-loop field probe.}
Freeze $\Lbase$ trajectories and evaluate all three frozen fields on identical states.  Compare
the final RL checkpoint with intermediate checkpoints and independent RL seeds.  A GenEval-perfect
checkpoint may be less transferable than an earlier checkpoint if it encodes reward shortcuts.

\paragraph{Stage 2: closed-loop zero-training oracle.}
For the deterministic diagnostic only, introduce a signed graft coefficient $\eta$ and sweep
\begin{equation}
\eta\in
\{-0.5,-0.25,0,0.125,0.25,0.5,1.0\}.
\label{eq:signed-graft-sweep}
\end{equation}
Use $\eta$ in place of $\lambda$ in \eqref{eq:inference-graft}.  Negative values are sign controls
for the ODE field and are outside the one-sided kernel-ratio theory of
Section~\ref{sec:kernel}; a negative kernel tilt would require mutual absolute continuity and
finite negative moments.
Include:
\begin{itemize}[leftmargin=1.6em]
\item a compute-matched $\Lbase$ run with more recipient-only NFEs;
\item a $\Sbase-\Sbase$ sham delta;
\item prompt-shuffled and norm-matched random deltas;
\item absolute $\Srl$ and $\Sbase$ fields scaled to the same RMS;
\item negative $\eta$, which should tend to reverse a genuine RL-improvement direction.
\end{itemize}

\paragraph{Stage 3: recipient training.}
Only if two neighboring positive strengths improve the reward--quality frontier, train
\eqref{eq:main-training-loss}.  Use the smallest strength within uncertainty of the best oracle
result and refresh recipient rollouts frequently.

\paragraph{Stage 4: alternative representations.}
If field grafting fails because donor queries are OOD on recipient states, test
flow-map CFM~\eqref{eq:flow-map-cfm}.  Attempt endpoint ratio/DMD only after verifying ratio
calibration and adequate ESS.

\subsection{Evaluation and failure modes}

For the GenEval example, report overall and per-category GenEval together with prompt alignment,
PickScore, OCR, aesthetics/realism, diversity, duplicate rate, and editing performance.  A perfect
GenEval score is not sufficient evidence of transferable general capability.

Likely failures include:
\begin{itemize}[leftmargin=1.6em]
\item the 9B RL delta encodes detector-specific shortcuts rather than compositional improvement;
\item FLUX.2-dev may already have little GenEval headroom, which must be measured under the exact
      evaluation protocol;
\item recipient trajectories leave the donor's reliable state region, especially late in
      denoising;
\item Qwen-conditioned donor corrections are not compatible with Mistral-conditioned recipient
      behavior;
\item the final RL checkpoint overshoots while an intermediate checkpoint transfers;
\item a numerically small delta is dominated by bf16 or nondeterministic-forward noise;
\item an open-loop-safe direction becomes unstable after closed-loop state drift;
\item a fixed global $\lambda$ spends its scale on timestep or resolution differences;
\item field loss improves while reward and perceptual quality degrade, demonstrating objective
      mismatch rather than implementation success.
\end{itemize}

\begin{tcolorbox}[colback=gray!5,colframe=gray!55,title=Recommended first implementation]
Use inference-time ODE velocity grafting as the causal oracle.  If it succeeds, distill the
recipient residual on recipient on-policy states.  Keep Euler-SDE or SDE-DPM-Solver++ as the
probabilistically exact transition arm and scale-calibration reference.  Do not begin with
endpoint density-ratio/DMD: it is more expensive and cannot repair an intrinsically
non-transferable donor shift.
\end{tcolorbox}

\subsection{Flow-Factory implementation}

The framework implementation uses trainer key
\texttt{flow-direct-opd} and lives under
\path{src/flow_factory/trainers/fdopd/}.  The canonical configuration is
\path{examples/flow_direct_opd/lora/flux2/default.yaml}.  It uses the Multi-Reward donor LoRA
\begin{center}
\small
\url{https://huggingface.co/Tencent-Hunyuan-Multimodal-RL/FLUX2-klein-base-9b-GenEval2-Multi-Reward}
\end{center}
by default; replacing only \texttt{donor\_rl\_lora\_path} with the corresponding Single-Reward
repository gives the first policy-shift ablation.

The implemented V1 dynamics are ODE velocity residuals and the existing Flow-SDE, Dance-SDE, and
CPS transition-mean wrappers.  DPM-SDE remains a theoretical extension because the current
Flow-Factory scheduler registry does not expose a DPM-SDE augmented-history transition.

%% ============================================================
\section{Conclusion}
\label{sec:conclusion}
%% ============================================================

The small RL checkpoint should be viewed as a cheap instrument that discovers a policy-improvement
direction, not as an absolute teacher for the larger model.  The general transferable signal is
the conditional kernel ratio $\dd K_{\Srl}/\dd K_{\Sbase}$.  Its geometric tilt of the large
reference kernel is an exact KL-regularized objective without any Gaussian assumption.  Shared
Gaussian SDE and DPM-SDE kernels make the target computationally simple: the donor transition-mean
shift is added to the large reference mean, and the resulting training loss is a covariance-weighted
residual MSE.  Continuous SDEs give the same local drift target through Girsanov, while the exact
global path tilt additionally contains a continuation-value correction.

For ODE flow models, transition ratios are not defined.  Nevertheless, the same velocity arithmetic
is a principled rescaled small-noise control limit and a valid deterministic field construction.
Flow-map CFM and endpoint density-ratio methods provide alternative deterministic and distributional
targets when the Eulerian donor shift is unreliable.  The practical research question is therefore
not whether a 9B model is a good teacher for a 32B model, but whether the \emph{change discovered by
9B RL} is stable and useful on the 32B model's own states.

\begin{thebibliography}{10}

\bibitem{directopd}
Shiyuan Feng, Huan-ang Gao, Haohan Chi, et al.
\newblock Weak-to-Strong Generalization via Direct On-Policy Distillation.
\newblock 2026.
\newblock \url{https://arxiv.org/abs/2607.05394}.

\bibitem{proxytuning}
Alisa Liu, Xiaochuang Han, Yizhong Wang, Yulia Tsvetkov, Yejin Choi, and Noah A. Smith.
\newblock Tuning Language Models by Proxy.
\newblock 2024.
\newblock \url{https://arxiv.org/abs/2401.08565}.

\bibitem{rast}
Siru Ouyang, Xinyu Zhu, Zilin Xiao, Minhao Jiang, Yu Meng, and Jiawei Han.
\newblock RAST: Reasoning Activation in LLMs via Small-model Transfer.
\newblock 2025.
\newblock \url{https://arxiv.org/abs/2506.15710}.

\bibitem{thinklogit}
Yunxiang Zhang, Muhammad Khalifa, Lechen Zhang, et al.
\newblock Logit Arithmetic Elicits Long Reasoning Capabilities Without Training.
\newblock \emph{Findings of ACL}, 2026.
\newblock \url{https://arxiv.org/abs/2507.12759}.

\bibitem{logitgrafting}
Apurv Verma, Binh-Nguyen Nguyen, Lingxiao Wang, and NhatHai Phan.
\newblock Logit Grafting: The Post-Training Delta is Sparse, Portable, and Powerful.
\newblock \emph{2nd Workshop on Compositional Learning at ICML}, 2026.
\newblock \url{https://openreview.net/forum?id=yWJLnCEKBI}.

\bibitem{flowmatching}
Yaron Lipman, Ricky T. Q. Chen, Heli Ben-Hamu, Maximilian Nickel, and Matt Le.
\newblock Flow Matching for Generative Modeling.
\newblock \emph{International Conference on Learning Representations}, 2023.
\newblock \url{https://arxiv.org/abs/2210.02747}.

\bibitem{dpmsolverpp}
Cheng Lu, Yuhao Zhou, Fan Bao, Jianfei Chen, Chongxuan Li, and Jun Zhu.
\newblock DPM-Solver++: Fast Solver for Guided Sampling of Diffusion Probabilistic Models.
\newblock 2022.
\newblock \url{https://arxiv.org/abs/2211.01095}.

\bibitem{dmd}
Tianwei Yin, Micha\"el Gharbi, Richard Zhang, Eli Shechtman, Fr\'edo Durand,
William T. Freeman, and Taesung Park.
\newblock One-step Diffusion with Distribution Matching Distillation.
\newblock \emph{CVPR}, 2024.
\newblock \url{https://arxiv.org/abs/2311.18828}.

\bibitem{diffusionopd}
Quanhao Li, Junqiu Yu, Kaixun Jiang, et al.
\newblock DiffusionOPD: A Unified Perspective of On-Policy Distillation in Diffusion Models.
\newblock 2026.
\newblock \url{https://arxiv.org/abs/2605.15055}.

\bibitem{gradualft}
Gudrun Thorkelsdottir and Arindam Banerjee.
\newblock
Gradual Fine-Tuning for Flow Matching Models.
\newblock 2026.
\newblock \url{https://arxiv.org/abs/2601.22495}.

\end{thebibliography}

\end{document}
